% !TEX program = lualatex
%% FullCourse_10Lecs -- Lecture 7, Act 2 (REPRESENT, part 1/2): Preprocessing & Bag-of-Words
%% NOTE: build with lualatex (renders the Chinese segmentation examples via luatexja+Fandol);
%% under pdflatex it still compiles, substituting the romanized fallbacks.
%% ch14 rebuild (2026-07-07), per Lec07_ch14_Rebuild_Plan_2026-07-07.md.
%% Reused frames copied verbatim (with renumbered cross-references) from
%%   archive_pre_split/pre_ch14_rebuild_2026-07-07/Lec07_T1a_Why_LLMs_Tokenization.tex (tokenization/preprocessing)
%%   archive_pre_split/pre_ch14_rebuild_2026-07-07/Lec07_T1b_Word_Embeddings.tex (TF/IDF/TF-IDF trio, relocated & condensed)
%% New frames follow chapter14 sec. 14.2-14.3.1: Chinese segmentation, DTM
%% sparsity/orthogonality, boundary + translate-vs-native thinkboxes, when-BoW-wins.

%% Shared preamble for the FullCourse_10Lecs topic decks.
%%
%% Each topic .tex \input's this file, then defines its own \title, \date
%% override (if any), \graphicspath, and \begin{document}.
%%
%% Reconstructed verbatim from the inlined copy preserved in
%% Lec10_Case_Studies/Lec10_T8_Case6_DeepRamsey.tex (lines 23-190), which is
%% explicitly annotated there as "from ../shared_preamble.tex, copied verbatim".
%% Base preamble matches MiniCourse_8hr/Slides/shared_preamble.tex; this version
%% additionally provides \sectiondivider, the conceptbox/cautionbox/examplebox/
%% takeawaybox tcolorboxes, and \compacttables used across the split decks.

\documentclass[10pt,english,aspectratio=169]{beamer}

\usepackage{ae,aecompl}
\usepackage[T1]{fontenc}
\usepackage{textcomp}
\usepackage[utf8]{inputenc}
\usepackage{babel}
\usepackage{datetime2}

\setcounter{secnumdepth}{3}
\setcounter{tocdepth}{3}

\usepackage{amsmath, amssymb, amsthm, graphicx}
\usepackage{booktabs, multirow}
\usepackage{tikz}
\usepackage{pgfplots}
\pgfplotsset{compat=1.16}
\usepackage[most]{tcolorbox}
\tcbuselibrary{skins, breakable, theorems}
\usepackage{listings}
\usepackage{xcolor}

\lstset{
  basicstyle=\ttfamily\small,
  keywordstyle=\color{blue!80!black}\bfseries,
  commentstyle=\color{green!50!black}\itshape,
  stringstyle=\color{orange!80!black},
  backgroundcolor=\color{gray!8},
  frame=single,
  rulecolor=\color{gray!40},
  breaklines=true,
  showstringspaces=false,
  numbers=none,
}

\lstdefinestyle{pythonstyle}{
    language=Python,
    basicstyle=\ttfamily\footnotesize,
    keywordstyle=\color{blue!80!black}\bfseries,
    commentstyle=\color{green!50!black}\itshape,
    stringstyle=\color{orange!80!black},
    frame=single,
    breaklines=true,
    showstringspaces=false,
    numbers=left,
    numbersep=5pt,
    numberstyle=\tiny\color{gray},
    backgroundcolor=\color{gray!8},
    rulecolor=\color{gray!40}
}

\ifx\hypersetup\undefined
  \AtBeginDocument{%
    \hypersetup{bookmarks=false,breaklinks=true,pdfborder={0 0 0},colorlinks=true,
      linkcolor=DarkText,urlcolor=RedTitle}%
  }
\else
  \hypersetup{bookmarks=false,breaklinks=true,pdfborder={0 0 0},colorlinks=true,
    linkcolor=DarkText,urlcolor=RedTitle}%
\fi

\makeatletter
\newcommand\makebeamertitle{%
  \begin{frame}[plain]
    \vspace*{1.0cm}
    \centering
    {\usebeamerfont{title}\usebeamercolor[fg]{title}\inserttitle\par}
    \vspace{1.0cm}
    {\usebeamerfont{author}\usebeamercolor[fg]{author}\insertauthor\par}
    \vspace{0.2cm}
    {\usebeamerfont{date}\usebeamercolor[fg]{date}\insertdate\par}
  \end{frame}}%
\AtBeginDocument{%
  \let\origtableofcontents=\tableofcontents
  \def\tableofcontents{\@ifnextchar[{\origtableofcontents}{\gobbletableofcontents}}
  \def\gobbletableofcontents#1{\origtableofcontents}%
}

\usetheme{metropolis}
\usefonttheme{professionalfonts}

\definecolor{LightGray}{RGB}{230,230,230}
\definecolor{RedTitle}{RGB}{0,0,200}
\definecolor{VeryLightGray}{RGB}{245,245,245}
\definecolor{DarkText}{RGB}{0,0,0}

\setbeamercolor{background canvas}{bg=white}
\setbeamercolor{title}{fg=RedTitle}
\setbeamercolor{author}{fg=DarkText}
\setbeamercolor{date}{fg=DarkText}
\setbeamercolor{frametitle}{bg=VeryLightGray,fg=DarkText}
\setbeamercolor{normal text}{fg=DarkText}
\setbeamerfont{title}{size=\large}

\setbeamersize{text margin left=5mm, text margin right=5mm}
\addtobeamertemplate{frametitle}{}{\vspace{-0.5em}}
\newcommand{\reducevspace}{\vspace{-0.5cm}}

% Shared section divider and box styles for split topic decks.
\newcommand{\sectiondivider}[2]{%
\begin{frame}[plain,noframenumbering]
\begin{center}
\vfill
{\Large\textbf{#1}\par}
\vspace{0.35cm}
{\normalsize #2\par}
\vfill
\end{center}
\end{frame}
}

\newtcolorbox{conceptbox}[1][]{
  colback=blue!3,
  colframe=RedTitle!55,
  boxrule=0.35mm,
  arc=1mm,
  left=2mm,
  right=2mm,
  top=1mm,
  bottom=1mm,
  #1
}
\newtcolorbox{cautionbox}[1][]{
  colback=red!3,
  colframe=red!50!black,
  boxrule=0.35mm,
  arc=1mm,
  left=2mm,
  right=2mm,
  top=1mm,
  bottom=1mm,
  #1
}
\newtcolorbox{examplebox}[1][]{
  colback=green!3,
  colframe=green!45!black,
  boxrule=0.35mm,
  arc=1mm,
  left=2mm,
  right=2mm,
  top=1mm,
  bottom=1mm,
  #1
}
\newtcolorbox{takeawaybox}[1][]{
  colback=VeryLightGray,
  colframe=RedTitle!55,
  boxrule=0.35mm,
  arc=1mm,
  left=2mm,
  right=2mm,
  top=1mm,
  bottom=1mm,
  #1
}
\newcommand{\compacttables}{\renewcommand{\arraystretch}{1.15}}

\setbeamertemplate{navigation symbols}{}
\setbeamertemplate{footline}{%
  \hfill%
  \usebeamercolor[fg]{page number in head/foot}%
  \usebeamerfont{page number in head/foot}%
  \insertframenumber\,/\,\inserttotalframenumber\kern1em\vskip2pt%
}

\makeatother

\author{Zhigang Feng}
% "date with time": compile timestamp so each build is identifiable.
\date{\today\ \DTMcurrenttime}


% --- CJK: real characters under Lua/XeLaTeX; romanized fallback under pdfLaTeX ---
\usepackage{iftex}
\ifluatex
  \usepackage{luatexja-fontspec}
  \setmainfont{Latin Modern Roman}%
  \setsansfont{Latin Modern Sans}%
  \setmonofont{Latin Modern Mono}%
  \setmainjfont{FandolSong-Regular.otf}[BoldFont=FandolHei-Regular.otf]%
  \newcommand{\zhterm}[2]{#1}%
\else\ifxetex
  \usepackage{fontspec}
  \setmainfont{Latin Modern Roman}%
  \setsansfont{Latin Modern Sans}%
  \setmonofont{Latin Modern Mono}%
  \usepackage{xeCJK}
  \setCJKmainfont{FandolSong-Regular.otf}[BoldFont=FandolHei-Regular.otf]%
  \newcommand{\zhterm}[2]{#1}%
\else
  \newcommand{\zhterm}[2]{#2}% pdflatex: romanized fallback
\fi\fi

\metroset{sectionpage=none}

\graphicspath{{./pic/}{../../MiniCourse_8hr/Slides/Module3_LLM_Text/pic/}{../}}

\usetikzlibrary{arrows.meta,positioning,shapes,calc,fit,backgrounds}

%% Lec07 shared MAP slide: "one map, five acts" recurring diagram.
%% Adapted from ../Lec09_Agentic_AI/lec09_map.tex (\LecNineMap -> \LecSevenMap).
%% Input this file AFTER ../shared_preamble.tex (needs RedTitle, VeryLightGray)
%% and call \LecSevenMap{k} inside a frame, where k in {1..5} highlights the
%% current act ("you are here") and k = 0 renders the closing reprise with all
%% five acts complete (used by T5b's closing frame).
%% Filename deliberately does not match the build glob Lec*_T*.tex, so the
%% build script never tries to compile it standalone.

\newcommand{\LecSevenMap}[1]{%
% Precompute per-box style selectors; TikZ option lists are comma-split before
% TeX conditionals run, so \ifnum must not sit inside a [...] options list.
\ifnum#1=0
  \tikzset{lecmapa/.style={lecmapdone}, lecmapb/.style={lecmapdone},
           lecmapc/.style={lecmapdone}, lecmapd/.style={lecmapdone},
           lecmape/.style={lecmapdone}}%
\else
  \tikzset{lecmapa/.style={}, lecmapb/.style={}, lecmapc/.style={},
           lecmapd/.style={}, lecmape/.style={}}%
  \ifnum#1=1 \tikzset{lecmapa/.style={lecmapcur}}\fi
  \ifnum#1=2 \tikzset{lecmapb/.style={lecmapcur}}\fi
  \ifnum#1=3 \tikzset{lecmapc/.style={lecmapcur}}\fi
  \ifnum#1=4 \tikzset{lecmapd/.style={lecmapcur}}\fi
  \ifnum#1=5 \tikzset{lecmape/.style={lecmapcur}}\fi
\fi
\begin{center}
\begin{tikzpicture}[
    act/.style={rectangle, rounded corners, draw=black!60, fill=VeryLightGray,
        minimum width=2.6cm, minimum height=0.92cm, text width=2.5cm,
        align=center, font=\scriptsize, text=DarkText},
    lecmapcur/.style={draw=RedTitle, very thick, fill=orange!20},
    lecmapdone/.style={draw=green!45!black, thick, fill=green!10},
    q/.style={font=\tiny\itshape, text width=2.7cm, align=center, text=black!70},
    sat/.style={rectangle, rounded corners, draw=black!45, dashed, fill=white,
        text width=5.0cm, align=center, font=\tiny, text=black!70,
        minimum height=0.5cm},
    barlab/.style={font=\tiny\bfseries, text=black!60},
    arr/.style={-{Stealth[length=2mm]}, thick, gray!60}
]
    % --- five act boxes ---
    \node[act, lecmapa] (a1) at (0,0)
        {\textbf{7.1 FRAME}\\ \tiny text as measurement};
    \node[act, lecmapb] (a2) at (2.85,0)
        {\textbf{7.2 REPRESENT}\\ \tiny tokens $\to$ vectors};
    \node[act, lecmapc] (a3) at (5.7,0)
        {\textbf{7.3 ATTEND}\\ \tiny attention \& Transformer};
    \node[act, lecmapd] (a4) at (8.55,0)
        {\textbf{7.4 PRETRAIN}\\ \tiny LMs \& alignment};
    \node[act, lecmape] (a5) at (11.4,0)
        {\textbf{7.5 MEASURE}\\ \tiny limits, apps, validation};
    \draw[arr] (a1) -- (a2);
    \draw[arr] (a2) -- (a3);
    \draw[arr] (a3) -- (a4);
    \draw[arr] (a4) -- (a5);
    % --- the student question each act answers ---
    \node[q] at (0,-0.82)    {what generates\\ this text?};
    \node[q] at (2.85,-0.82) {how do we turn it\\ into numbers?};
    \node[q] at (5.7,-0.82)  {how does context\\ get encoded?};
    \node[q] at (8.55,-0.82) {where does a pretrained\\ model come from?};
    \node[q] at (11.4,-0.82) {is my measure valid ---\\ and defensible?};
    % --- "you are here" marker above the current act ---
    \ifnum#1=1 \node[font=\scriptsize\bfseries, text=RedTitle] at (0,0.95)
        {you are here\ $\blacktriangledown$};\fi
    \ifnum#1=2 \node[font=\scriptsize\bfseries, text=RedTitle] at (2.85,0.95)
        {you are here\ $\blacktriangledown$};\fi
    \ifnum#1=3 \node[font=\scriptsize\bfseries, text=RedTitle] at (5.7,0.95)
        {you are here\ $\blacktriangledown$};\fi
    \ifnum#1=4 \node[font=\scriptsize\bfseries, text=RedTitle] at (8.55,0.95)
        {you are here\ $\blacktriangledown$};\fi
    \ifnum#1=5 \node[font=\scriptsize\bfseries, text=RedTitle] at (11.4,0.95)
        {you are here\ $\blacktriangledown$};\fi
    \ifnum#1=0 \node[font=\scriptsize\bfseries, text=green!45!black] at (5.7,0.95)
        {all five acts complete};\fi
    % --- bottom band: the measurement sandwich ---
    \fill[blue!10]   (-1.3,-1.55) rectangle (1.40,-1.22);
    \fill[green!10]  (1.65,-1.55) rectangle (9.85,-1.22);
    \fill[orange!15] (10.1,-1.55) rectangle (12.7,-1.22);
    \node[barlab] at (0.05,-1.385) {WHY TEXT = MEASUREMENT};
    \node[barlab] at (5.75,-1.385) {HOW THE MACHINE READS TEXT};
    \node[barlab] at (11.4,-1.385) {MEASURE, VALIDATE, DISCLOSE};
    % --- the recurring spine (text DGP), printed under the band ---
    \node[font=\tiny\itshape, text=black!60] at (5.7,-1.83)
        {the spine:\ \ $x_i = g(s_i,\,a_i,\,c_i)+\varepsilon_i
         \;\longrightarrow\; m(x_i)
         \;\longrightarrow\; y_i = \beta\,m(x_i)+\gamma z_i+u_i$};
    % --- dashed satellites: the neighbouring lectures ---
    \node[sat] at (2.0,-2.42)
        {\textit{before 7.1:} Lec03--04 gave you neural nets, backprop \& PyTorch --- Lec07 teaches them to read};
    \node[sat] at (9.8,-2.42)
        {\textit{after 7.5:} Lec08 RAG (grounded memory) $\to$ Lec09 agents (grounded action) $\to$ Lec10 cases};
\end{tikzpicture}
\end{center}
}


\title[AI for Econ Research]{AI for Economic Research: Dynamic Models, Language, and Agents\\
Lecture 7.2a: Preprocessing and Bag-of-Words}

\begin{document}

\makebeamertitle

%=============================================================================
\begin{frame}{Lecture 7 in One Map}
\LecSevenMap{2}
\medskip
\textit{\footnotesize Route: raw strings $\to$ clean text $\to$ tokens $\to$ counts (part 1/2 of REPRESENT --- this deck builds the transparent count-based representations; part 2/2 turns counts into dense vectors).}
\end{frame}

%=============================================================================
\section{From Strings to Tokens}
\sectiondivider{From Strings to Tokens}{Tokenization, subwords --- and the Chinese special case}

\begin{frame}{Tokenization: The First Step}
\begin{itemize}
\item \textbf{Goal:} convert raw character sequences into discrete \emph{tokens} that a neural network can consume.

\medskip
\item \textbf{Why not just split by spaces?}
\begin{itemize}
    \item Many economically relevant terms combine words, numbers, and symbols (e.g.\ \texttt{S\&P500}, \texttt{CO2}, \texttt{Q4\_2025}, \texttt{stagflation})
    \item Punctuation and formatting carry meaning (`\$1.2B' vs `1.2 billion')
    \item Different languages have different word boundaries (or none at all)
\end{itemize}

\medskip
\item \textbf{What is a token?} A flexible unit:
\begin{itemize}
    \item A whole word (`inflation')
    \item A subword (`infl', `ation')
    \item A single character or byte
    \item A special marker (\texttt{<EOS>}, \texttt{<PAD>}, \texttt{<USER>})
\end{itemize}
\end{itemize}
\end{frame}

\begin{frame}{Vocabulary and the Tokenization Map}
\begin{itemize}
\item Let $A$ be the alphabet (Unicode characters). A tokenizer is a function:
\[
   \mathcal{K}\colon A^{*} \;\longrightarrow\; \{1, 2, \ldots, M\}^{*}
\]
mapping any text string into a sequence of \textbf{numbers} --- each number is a positive integer that serves as an index into a lookup table of size $M$. The fact that the index is an integer is incidental; what matters is that it \emph{uniquely identifies} each token, the same way a ZIP code uniquely identifies a location.

\medskip
\item \textbf{Why integers?} Neural networks need numeric inputs; each token ID will index into an embedding lookup table.

\medskip
\item \textbf{Practical constraint --- context length:}
\begin{itemize}
    \item Models have a maximum input size $n$ (typically $2{,}048$ to $128{,}000$ tokens; some $1$M)
    \item Long policy documents, FOMC transcripts, or 10-K filings often exceed $n$
    \item Forces \emph{chunking}, \emph{summarization}, or \emph{retrieval} (a hint at Lec08)
\end{itemize}

\medskip
\item Vocabulary sizes are typically $30{,}000$--$100{,}000$ subword tokens. \textbf{To put this in perspective:} Merriam-Webster Unabridged lists $\approx470{,}000$ entries; a college-educated native speaker actively uses roughly 20{,}000--40{,}000 words. A subword vocabulary of $M=50{,}000$ therefore spans roughly the active vocabulary of a well-read person --- but because it tokenizes \emph{sub}words (stems, suffixes, common fragments), it effectively covers far more word forms and technical terms than a same-size word-level vocabulary would.
\end{itemize}
\end{frame}

\begin{frame}{Illustrative Example: ``GDP rose by 2\%''}
\begin{itemize}
\item A simple tokenization might produce
\[
   \text{(GDP},\; \text{rose},\; \text{by},\; \text{2},\; \%\text{)}
\]
mapping to integer IDs, e.g.\ $(541,\; 823,\; 71,\; 312,\; 49)$.

\medskip
\item \textbf{Why integer mapping matters:}
\begin{itemize}
    \item Each ID becomes the index of a learned embedding vector
    \item Numeric tokens (`2', `3', `5') get \emph{their own} embeddings
\end{itemize}

\medskip
\item \textbf{Injectivity matters for forecasting:} a well-designed tokenizer maps
\begin{itemize}
    \item ``inflation rate is 3.5\%''
    \item ``inflation rate is 3.6\%''
\end{itemize}
to \emph{different} token sequences. Otherwise small but economically meaningful differences would collapse into the same input.

\medskip
\item \textbf{Tokenizer/model coupling:} the tokenizer is part of the model. If you swap the tokenizer, all the embedding rows are meaningless --- this is why every published LLM ships with its tokenizer.
\end{itemize}
\end{frame}

\begin{frame}{Chinese Word Segmentation: Boundaries Must Be Inferred}
\small
\textbf{English gets word boundaries for free; Chinese does not --- so segmentation is itself a model.}
\begin{itemize}
\item Chinese runs characters together; word boundaries must be \emph{inferred}. Segmentation is a sequence-labeling problem: character by character, where does a word end?
\item The classic ambiguity --- \zhterm{研究生命科学}{yanjiu-sheng-ming-ke-xue} splits as
\begin{itemize}
    \item \zhterm{研究\,$\vert$\,生命科学}{yanjiu $\vert$ shengming-kexue} --- ``research \emph{on} life science'', or
    \item \zhterm{研究生\,$\vert$\,命科学}{yanjiusheng $\vert$ ming-kexue} --- ``graduate student'' $+$ a nonsense remainder.
\end{itemize}
Economic text has its own: \zhterm{行业政策支持}{hangye zhengce zhichi} parses as \zhterm{行业政策\,$\vert$\,支持}{hangye-zhengce $\vert$ zhichi} (``industry-policy'' $+$ ``support'') or \zhterm{行业\,$\vert$\,政策支持}{hangye $\vert$ zhengce-zhichi} (``industry'' $+$ ``policy-support'') --- different head word, different meaning.
\item \textbf{Tools:} Jieba, THULAC, LTP --- each a dictionary plus a statistical/neural sequence labeler, trading speed against accuracy.
\item \textbf{The economist's trap --- domain compounds:} terms like ``monetary policy committee'', ``macro-prudential framework'', ``aggregate financing to the real economy'' (\zhterm{社会融资规模}{shehui rongzi guimo}) are not in general-purpose dictionaries; one mis-split changes the head word.
\item \textbf{Fix:} load your domain glossary into the tool's \textbf{user dictionary} so policy terms survive as whole tokens. Stopwords differ too: drop structural particles (\zhterm{的}{de}, \zhterm{了}{le}, \zhterm{是}{shi}); keep logical connectives (\zhterm{因此}{yinci} ``therefore'', \zhterm{然而}{ran'er} ``however'') for information-extraction tasks.
\end{itemize}
\end{frame}

%=============================================================================
\section{Preprocessing Is a Modeling Decision}
\sectiondivider{Preprocessing Is a Modeling Decision}{Every cleaning switch changes the downstream estimate}

\begin{frame}{Tokenization Is a Preprocessing Choice}
\begin{itemize}
\item Tokenization sits inside a larger \textbf{preprocessing pipeline}:
\[
   x^{\text{raw}} \;\xrightarrow{\;\text{clean}\;}\; x^{\text{clean}}
   \;\xrightarrow{\;\text{tokenize}\;}\; \{t_{ij}\}
   \;\xrightarrow{\;\text{represent}\;}\; z_i,
\]
from raw scraped text $\to$ normalized text $\to$ integer token IDs $\to$ downstream numeric representation (TF-IDF, embedding, $\ldots$).

\end{itemize}
\end{frame}

\begin{frame}{The Seven Preprocessing Choices}
\small
\begin{itemize}
\item \textbf{Lowercasing} --- map \texttt{Fed} and \texttt{fed} to the same token. Compresses vocabulary, but erases proper-noun signal.
\item \textbf{Stemming / lemmatization} --- cut suffixes so \emph{banking/banks/banked} collapse to \emph{bank}. Can erase tense, plurality, policy-language nuance.
\item \textbf{Stop-word removal} --- drop high-frequency function words. Reduces noise; but negation (\emph{not}) and modality (\emph{may, will}) can matter for central-bank text.
\item \textbf{Punctuation handling} --- keep, strip, or tokenize separately. Percent signs, decimals, sentence boundaries all matter.
\item \textbf{$n$-gram inclusion} --- treat \emph{``yield curve''} as one token rather than two. Captures multi-word economic concepts at the cost of a larger vocabulary.
\item \textbf{Frequency cutoffs} --- drop terms in $<k$ or $>K$ documents. Removes typos and boilerplate; can also remove rare-but-informative events.
\item \textbf{Language filtering} --- restrict to English or apply per-language pipelines. Crucial for multilingual corpora like ECB speeches.
\end{itemize}

\end{frame}

\begin{frame}{One Token Boundary Can Flip the Sign}
\textbf{Where you draw token boundaries can reverse an economic conclusion.}
\begin{cautionbox}
FOMC-style sentence: \emph{``\ldots growth is \textbf{not sufficient} to warrant further tightening.''}
\begin{itemize}
\item \textbf{Strategy A (unigrams):} \texttt{not} and \texttt{sufficient} are separate features. If \texttt{sufficient} loads positively in training, a bag-of-words model can extract a \emph{positive} signal from a negation.
\item \textbf{Strategy B (bigrams):} \texttt{not\_sufficient} is one feature --- the negation survives as a unit.
\end{itemize}
\end{cautionbox}
\medskip
\begin{itemize}
\item Does the choice still matter under BERT-class models, whose attention can connect \texttt{not} to \texttt{sufficient} across positions? (Preview: mostly no --- one reason Acts 3--4 exist.)
\item Either way: the tokenizer version and configuration belong in the replication package --- they are part of the measurement instrument.
\end{itemize}
\end{frame}

%=============================================================================
\section{Bag-of-Words and TF-IDF}
\sectiondivider{Bag-of-Words and TF-IDF}{The transparent baseline: count, weight, audit}

\begin{frame}{Bag-of-Words: The Document--Term Matrix}
\textbf{The simplest representation: count words, forget order.}
\begin{itemize}
\item Toy corpus --- doc 1: \emph{``central bank cuts rates to support growth''}; doc 2: \emph{``central bank raises rates to curb inflation''}.
\end{itemize}
\begin{center}
\footnotesize
\begin{tabular}{l c c c c c c c c c}
\toprule
 & central & bank & cuts & raises & rates & support & curb & growth & inflation \\
\midrule
doc 1 & 1 & 1 & 1 & 0 & 1 & 1 & 0 & 1 & 0 \\
doc 2 & 1 & 1 & 0 & 1 & 1 & 0 & 1 & 0 & 1 \\
\bottomrule
\end{tabular}
\end{center}
\begin{itemize}
\item In general: $D$ documents $\times$ vocabulary of size $V$ (typically $V \gg 10^4$); entry $X_{d,v}$ is binary, a raw count, or a TF-IDF weight (two slides ahead).
\item \textbf{The assumption you buy:} ``central bank cuts rates to support growth'' $=$ ``growth supports central bank rate cuts''. Harmless when counting topic words (EPU); fatal for direction, causation, negation.
\end{itemize}
\end{frame}

\begin{frame}{Bag-of-Words: Two Built-in Consequences}
\textbf{Counting words and forgetting order leaves two statistical fingerprints on the matrix.}
\begin{itemize}
\item \textbf{Sparsity:} a document uses only 1--2\% of the vocabulary $\Rightarrow$ the matrix is $>$99\% zeros; with $D \ll V$, OLS on raw counts is not even identified --- dimension reduction (PCA, topic models) or regularization (Lasso) is mandatory.
\item \textbf{Orthogonality:} every word is its own perpendicular axis. \emph{Unemployment} $\perp$ \emph{joblessness}: two articles about the same phenomenon in different vocabulary register as \emph{unrelated}.
\end{itemize}
\medskip
\textit{\footnotesize Orthogonality is the recall problem behind every dictionary method --- and the motivation for the dense embeddings of T2b.}
\end{frame}

\begin{frame}{Term Frequency and Inverse Document Frequency}
\textbf{Weight words by how much they distinguish a document --- not by how often they appear.}
\begin{columns}[T]
\column{0.46\textwidth}
\textbf{TF} --- share of document $d$ that is word $w$:
\[ \mathrm{TF}(w,d) = \frac{\mathrm{count}(w,d)}{\text{total words in } d} \]
{\footnotesize \emph{``the Fed raised rates because inflation persisted''} (7 words): every word gets TF $=1/7\approx 0.14$ --- \texttt{the} scores as high as \texttt{inflation}.\\[3pt]
\textbf{Limitation:} TF alone rewards frequency, not information. We need a counterweight.}
\column{0.50\textwidth}
\textbf{IDF} --- rarity of $w$ across $N$ documents:
\[ \mathrm{IDF}(w) = \log\!\left(\frac{N}{\#\{\text{docs containing } w\}}\right) \]
{\footnotesize Corpus of $N=100$ FOMC speeches:
\begin{center}
\begin{tabular}{l c c}
\toprule
word & docs & IDF \\
\midrule
\emph{the} & 100 & $0$ \\
\emph{inflation} & 60 & $0.51$ \\
\emph{yield curve} & 5 & $3.00$ \\
\bottomrule
\end{tabular}
\end{center}
Rare $\Rightarrow$ informative $\Rightarrow$ up-weighted. \emph{the} zeroes out automatically --- no stop-word list needed.}
\end{columns}
\end{frame}

\begin{frame}{TF-IDF --- and When Counting Still Wins}
\textbf{Multiply term frequency by inverse document frequency; each document becomes a sparse vector in $\mathbb{R}^{|V|}$.}
\[
   \text{TF-IDF}(w, d) \;=\; \text{TF}(w, d)\,\cdot\,\text{IDF}(w)
\]
\smallskip
Orthogonality remains a real defect --- yet counts win whole classes of research problems:
\begin{itemize}
\item \textbf{When the target is ``how often do well-defined words appear'':} EPU (Thread A), the Loughran--McDonald finance dictionary, Hassan et al.'s bigram political risk --- all count-based \emph{by design}.
\item \textbf{The killer feature is auditability:} you can list exactly which article and which phrase triggered the index --- line-by-line external validation. No embedding offers that evidence chain.
\item \textbf{Also the sensible default when:} the corpus is small ($D < 1{,}000$); compute is tight; or the measure must plug into a standard econometric pipeline with transparent features.
\end{itemize}
\smallskip
\begin{takeawaybox}
Choose the representation for the \emph{role} the text plays (T1's three roles): measurement wants transparency first; behavior and nuance want context --- exactly where T2b and Act 3 head next.
\end{takeawaybox}
\end{frame}

%=============================================================================
\begin{frame}{Bridge: Counts Cannot See Synonyms}
\textbf{The bag of words has taken us as far as orthogonal axes can go.}
\begin{itemize}
\item In the document--term matrix, \emph{hawkish} and \emph{tightening} are perpendicular; \emph{unemployment} and \emph{joblessness} share nothing. Every count-based measure inherits this blindness --- the recall problem behind all dictionary methods.
\item \textbf{Next deck:} learn coordinates \emph{from co-occurrence}, so related words end up near each other --- Word2Vec, GloVe, and the geometry that lets us do arithmetic with meaning.
\item The DGP lens carries over unchanged: an embedding is just a richer $m(\cdot)$; the hazards $a_i$, $c_i$ do not vanish because the vectors get denser.
\end{itemize}
\bigskip
\centering\footnotesize\textit{Arc: T1 FRAME $\to$ \textbf{T2 REPRESENT} $\to$ T3 ATTEND $\to$ T4 PRETRAIN $\to$ T5 MEASURE.}
\end{frame}

\end{document}
